% BHU PYQ -- Manually transcribed & error-corrected
% Paper: MTB-609 : Special Theory of Relativity-II
% Exam: B.A./B.Sc. (Hons) Semester VI Examination, 2013-14

\documentclass[11pt,a4paper]{article}
\usepackage[a4paper,top=2.5cm,bottom=2.5cm,left=2.8cm,right=2.8cm,headheight=14pt]{geometry}
\usepackage{amsmath,amssymb}
\usepackage[T1]{fontenc}\usepackage[utf8]{inputenc}
\usepackage{lmodern,microtype,fancyhdr,enumitem,hyperref,lastpage,parskip}

\hypersetup{colorlinks=true,urlcolor=black,linkcolor=black,
  pdftitle={MTB-609: Special Theory of Relativity-II -- BHU B.A./B.Sc. Sem VI 2013-14}}
\pagestyle{fancy}\fancyhf{}
\renewcommand{\headrulewidth}{0.4pt}\renewcommand{\footrulewidth}{0.4pt}
\fancyhead[L]{\small\textit{Banaras Hindu University}}
\fancyhead[R]{\small\textit{MTB-609: Relativity-II}}
\fancyfoot[C]{\small Page \thepage\ of \pageref{LastPage}}

\newlist{parts}{enumerate}{2}
\setlist[parts,1]{label=(\alph*),leftmargin=*,itemsep=5pt,topsep=3pt}
\setlist[parts,2]{label=(\roman*),leftmargin=*,itemsep=3pt,topsep=2pt}
\newcommand{\pts}[1]{\hfill\textit{[#1]}}

\begin{document}
\begin{center}
  {\large\bfseries BANARAS HINDU UNIVERSITY}\\[2pt]
  {\normalsize B.A./B.Sc. (Hons) --- Semester VI Examination, 2013-14}\\[6pt]
  \rule{0.8\linewidth}{0.5pt}\\[6pt]
  {\large\bfseries Paper No. MTB-609: Special Theory of Relativity-II}\\[4pt]
  {\normalsize Time: 3 Hours \hfill Full Marks: 70 \quad (Symbols have their usual meanings)}
\end{center}
\rule{\linewidth}{0.4pt}
\smallskip
\noindent\textit{Instructions: Answer \textbf{five} questions including Question~1 which is compulsory. Figures in right-hand margin indicate marks.}
\bigskip

\noindent\textbf{Question 1.}
\begin{parts}
  \item A proton of rest mass $1.67 \times 10^{-24}\,\text{g}$ is moving with velocity $0.995c$. Calculate its energy and momentum. \pts{4}
  \item Show that the momentum $p$ of a particle of rest mass $m_0$ and kinetic energy $K$ is given by:
  \[
  p = \sqrt{\frac{K^2}{c^2} + 2m_0 K}
  \] \pts{4}
  \item Obtain the transformation equations for $\vec{H}$ (the magnetic field). \pts{4}
  \item Prove that $\vec{E}\cdot\vec{H}$ is Lorentz invariant. \pts{4}
  \item Find $F_{12}$ and $F_{13}$ in terms of $\vec{E}$ and $\vec{H}$. \pts{4}
  \item State the expression for the Lorentz force density. \pts{2}
\end{parts}

\medskip\rule{\linewidth}{0.2pt}

\medskip\noindent\textbf{Question 2.}
\begin{parts}
  \item Establish the transformation formula for relativistic mass and hence prove that $F = \dfrac{dp}{dt}$. \pts{6}
  \item Two particles with proper masses $m_1$, $m_2$ move along the $x$-axis with velocities $u_1$, $u_2$. They collide and coalesce. Using conservation of relativistic momentum and energy, prove that the proper mass $m_3$ and velocity $u_3$ of the resulting particle satisfy:
  \[
  m_3 = \sqrt{m_1^2 + m_2^2 + 2m_1 m_2 \gamma_1 \gamma_2\!\left(1 - \tfrac{u_1 u_2}{c^2}\right)}, \quad u_3 = \frac{m_1 \gamma_1 u_1 + m_2 \gamma_2 u_2}{m_1 \gamma_1 + m_2 \gamma_2}
  \]
  where $\gamma_i = \left(1 - \dfrac{u_i^2}{c^2}\right)^{-1/2}$. \pts{6}
\end{parts}

\medskip\rule{\linewidth}{0.2pt}

\medskip\noindent\textbf{Question 3.}
\begin{parts}
  \item Establish the transformation formulae for energy and momentum. \pts{6}
  \item What do you understand by the relativistic Hamiltonian? \pts{6}
\end{parts}

\medskip\rule{\linewidth}{0.2pt}

\medskip\noindent\textbf{Question 4.}
\begin{parts}
  \item Obtain the transformation formulae for the densities of electric charge and electric current. \pts{6}
  \item Prove the Lorentz invariance of Maxwell's equations without using tensors. \pts{6}
\end{parts}

\medskip\rule{\linewidth}{0.2pt}

\medskip\noindent\textbf{Question 5.}
\begin{parts}
  \item Obtain the transformation equations for $\vec{A}$ and $\phi$, where $\vec{A}$ is the electromagnetic vector potential and $\phi$ is the scalar potential. \pts{6}
  \item Find all components of the anti-symmetric tensor $F^{\mu\nu}$ in four dimensions, in terms of $\vec{E}$ and $\vec{H}$. \pts{6}
\end{parts}

\medskip\rule{\linewidth}{0.2pt}

\medskip\noindent\textbf{Question 6.}
Obtain the expression for the electromagnetic energy-momentum tensor in Minkowski space. \pts{12}

\medskip\rule{\linewidth}{0.2pt}

\medskip\noindent\textbf{Question 7.}
\begin{parts}
  \item Write a short note on the equivalence of mass and energy ($E = mc^2$). \pts{6}
  \item Using the Lorentz condition, show that $\Box^2 \phi = 4\pi\rho$ and $\Box^2 \vec{A} = \dfrac{4\pi}{c}\vec{J}$. \pts{6}
\end{parts}

\vfill\rule{\linewidth}{0.4pt}
\begin{center}\small
  \href{https://bhu-exam.vercel.app/}{Click here to visit our website}\\[4pt]
  \href{https://me-aryan.vercel.app/}{Click here to visit our contributor}\\[4pt]
  \href{https://quashan-me.vercel.app/}{Click here to visit our contributor}
\end{center}
\end{document}
